\section{Problem setting and evaluation protocol}
\label{sec:setup}

\paragraph{Setting.}
For a benchmark with $N$ items, let $M\in\{0,1\}^{K\times N}$ hold the correctness of $K$ models on every item. The models are partitioned into \emph{source} models $\mathcal S$, whose full rows are known, and \emph{target} models $\mathcal T$, which may be run on only $n\ll N$ items. A \emph{selector} chooses an item set $I$ with $|I|=n$ using $M_{\mathcal S}$ only (adaptive selectors may additionally use the target's answers already observed); a \emph{predictor} maps the target's outcomes on $I$ and $M_{\mathcal S}$ to an estimate $\hat y_t$ of the target's full accuracy $y_t=\frac1N\sum_i M_{ti}$. Target full scores are used only to score predictions; no selection, fitting or tuning step sees them. Hyperparameters are either fixed a priori (the IRT, low-rank and embedding sizes) or chosen on source models only (the ridge penalty and the design of P3, by cross-validation across source models); none is tuned on targets.

\paragraph{Metrics.}
For each benchmark, budget $n$ and split we report, over the target models, the mean absolute error (MAE) and root mean squared error (RMSE) of $\hat y_t$ against $y_t$ in percentage points of accuracy, and the Spearman and Kendall correlations between $\hat y$ and $y$.

\paragraph{Data.}
We use the item-level 0/1 score matrices of the benchmark-prediction repository of \citet{zhang2025benchmark} at the pinned commit \texttt{55faff4}, covering seven Open LLM Leaderboard v2 benchmarks and six HELM benchmarks (Table~\ref{tab:data}); model accuracies span a wide range, but are narrow for GPQA, MuSR and MATH Lvl5. ``Constant items'' are items on which all models agree; they are kept in $M$ and in $y_t$.

\begin{table}[t]
\caption{Score matrices used (model accuracies in \%). Generated by \texttt{figures/make\_data\_table.py}.}
\label{tab:data}
\centering\small
\begin{tabular}{lrrcr}
\toprule
Benchmark & Models & Items & Model accuracy (min / median / max) & Constant items \\
\midrule
IFEval & 448 & 541 & 3.7 / 27.8 / 88.0 & 1.3\% \\
BBH & 446 & 5761 & 24.1 / 44.2 / 74.7 & 0.0\% \\
GPQA & 447 & 1192 & 22.5 / 28.2 / 43.7 & 0.0\% \\
MuSR & 448 & 756 & 32.3 / 40.0 / 53.6 & 1.3\% \\
MATH Lvl5 & 447 & 894 & 0.0 / 2.1 / 35.6 & 13.3\% \\
MMLU-Pro & 448 & 12032 & 10.3 / 27.8 / 59.7 & 0.0\% \\
ARC-C & 218 & 1172 & 20.5 / 56.4 / 75.1 & 2.7\% \\
\midrule
HELM MMLU & 83 & 567 & 22.9 / 63.8 / 81.1 & 0.2\% \\
HELM GSM8K & 83 & 1000 & 2.8 / 73.4 / 95.6 & 0.3\% \\
HELM MATH & 83 & 437 & 2.5 / 65.0 / 93.4 & 0.2\% \\
HELM MedQA & 83 & 1000 & 21.5 / 61.0 / 86.9 & 0.2\% \\
HELM LegalBench & 79 & 2047 & 22.8 / 56.8 / 70.2 & 2.6\% \\
HELM Commonsense & 83 & 500 & 22.2 / 87.8 / 97.2 & 0.0\% \\
\bottomrule
\end{tabular}

\end{table}
